%% kapa: exam=1 round=34 year=2023 period=1 form=A subject=civil_law subject_ko=민법 seq=4 number=4
%% source: Q-Net 감정평가사 기출문제 (cst00309) artlSeq=5211789 file=raw/1008/5211789/2023년도 제34회 감정평가사 1차시험 1교시.pdf
%% conversion: kordoc 4.24.51 -> tools/gen1.py (auto); review=proofread-sonnet; flags=none
\begin{question}{4}
민법상 법인의 정관에 관한 설명으로 옳은 것은? (다툼이 있으면 판례에 따름)
\begin{choices}
  \item 감사의 임면에 관한 사항은 정관의 필요적 기재사항이다.
  \item 정관의 임의적 기재사항은 정관에 기재되더라도 정관의 변경절차 없이 변경할 수 있다.
  \item 정관변경의 의결정족수가 충족되면 주무관청의 허가가 없어도 정관변경의 효력이 생긴다.
  \item 재단법인이 기본재산을 편입하는 행위는 주무관청의 허가를 받지 않아도 유효하다.
  \item 재단법인의 기본재산에 관한 저당권 설정행위는 특별한 사정이 없는 한 주무관청의 허가를 얻을 필요가 없다.
\end{choices}
\end{question}
\answer{5}
